\scanpage{240}
\mspage{186}
% archive sheet number: 239
jusqu'à aujourd'hui même, car j'avais l'idée
que, grâce au mirifique foncteur $C$, la \unsure{catégori|catégorie}
$\underline{\mathrm{Parf}}_W\,S$, comme les sept autres catégories de l'étage,
pouvait servir à définir $\mathrm{HOT}_W^{\mathrm{lc}}(S)$. Or il n'en est
rien, semble-t-il, à moins de \deleted{\unsure{faire}}\inserted{supposer} l'axiome
$W(4\mathrm b)$, ou \deleted{\unsure{bien}} variante $W(4\mathrm b')$ (nouvelle
version). L'ancienne version de $W(4\mathrm b')$ \deleted{\unsure{était}}\inserted{s'avère}
parfaitement inutilisable.

A ce propos, je n'ai tjrs pas répondu clairement
pourtant, par un ``oui'' ou un ``non'', à la question :

\uline{Question} : les inclusions
\diagram{\begin{tikzcd}[ampersand replacement=\&,column sep=small,row sep=small]
 \&\underline{\mathrm{Liss}}_{0W}\,S\\
 \underline{\mathrm{Parf}}_W\,S\arrow[ur]\arrow[dr]\&\\
 \&\underline{\mathrm{Prop}}_{0W}\,S
\end{tikzcd}}
induisent-elles, même en l'absence de $W(4\mathrm b)$
(ou de ses variantes affaiblies) des équivalences
pour les localisés fins (par $W_S^f$) ? Même
question pour $\underline{\mathrm{Parf}}_{\unsure{\infty|W}}\,S$ au lieu de $\underline{\mathrm{Parf}}_W\,S$.* 
% slanted
\margin{OK en présence de l'axiome $W(6\mathrm{bis})$ en ce \uline{\unsure{lieu|sens}}, cf p. 192}
Même question pour
\[
\underline{\mathrm{Bifib}}\,S\subset\underline{\mathrm{Parf}}_W\,S,
\]
la catégorie des catégories \deleted{\unsure{S}} sur $S$ qui sont à
la fois $\mathrm{Cat}$-fibrées et $\mathrm{Cat}$-cofibrées. \deleted{\unsure{Si}}

Je ne me rends pas compte, \deleted{si les} même en
admettant (cette fois) l'axiome $W(4\mathrm a)$, si
$\underline{\mathrm{Bifib}}\,S$ est assez grosse pour pouvoir décrire tous
les types d'homotopie \inserted{relatifs} loc. const. sur $S$, i.e. si
\[
(\underline{\mathrm{Bifib}}\,S)(W_S^f)^{-1}\mathrel{\unsure{\text{---}|\longrightarrow}}
(\underline{\mathrm{Fib}}_W\,S)(W_S^f)^{-1}
\]
est au moins ess. surjectif, sinon une équivalence de
catégories. Il manque un foncteur
\[
\mathrm{Cat}_{/S}\mathrel{\unsure{\text{---}|\longrightarrow}}\underline{\mathrm{Bifib}}\,S
\qquad\text{ou du moins}\qquad
\underset{\inserted{\substack{\text{ou}\\\underline{\unsure{\mathrm{Fis}|\mathrm{Fib}}}_W\,S\,?}}}{\underset{\llcorner\!\rule{2em}{.35pt}\!\lrcorner}{\underline{\mathrm{Fib}}_W\,S}}
\mathrel{\unsure{\text{---}|\longrightarrow}}\underline{\mathrm{Bifib}}\,S,
\]
qui \deleted{\unsure{joue}} puisse jouer le rôle de $\Phi$ et de $\Psi$ \uline{simultané}%
